
\section{Cinemática inversa}
\label{sec:cin_inv}

A função \texttt{inverse\_kin} implementa a formulação matemática da cinemática inversa apresentada no Capítulo~\ref{ch:inv_kin}, permitindo determinar os ângulos articulares necessários para alcançar uma posição pretendida do centro de preensão.

A resolução é realizada com o auxílio das funções \texttt{ik\_2dof} e \texttt{solve\_pitch}. A primeira resolve o problema geométrico de um manipulador planar de dois graus de liberdade, enquanto a segunda determina os ângulos das juntas J1 a J4 para uma inclinação específica da garra, verificando as restrições de alcance e os limites articulares.

A função principal começa por calcular a distância frontal e a altura relativamente ao ombro. Seguidamente, procura uma inclinação válida para a garra, começando pela orientação preferencial de $90^\circ$ e testando, quando necessário, inclinações alternativas em incrementos de $5^\circ$.

Quando é encontrada uma solução válida numa inclinação diferente da preferencial, é realizado um refinamento através do método da bisseção. Por fim, é calculada a rotação da junta J6, caso tenha sido fornecida a orientação do objeto e a inclinação da garra satisfaça a condição estabelecida.

A implementação da função principal encontra-se apresentada na Listagem~\ref{lst:inv_kin}.

\begin{lstlisting}[
    language=Matlab,
    caption={Implementação da função inverse\_kin.},
    label={lst:inv_kin}
]
function [armJoints, isReachable, pitchUsed] = ...
    inverse_kin(x, y, z, Links, qMin, qMax, ...
                objYaw, prefPitch, family, fixedPitch)

    % Valores por defeito
    if nargin < 5 || isempty(qMin)
        qMin = deg2rad([-165 -165 -165 -165 -165 -175]);
    end
    if nargin < 6 || isempty(qMax)
        qMax = deg2rad([165 165 165 165 165 175]);
    end
    if nargin < 7
        objYaw = [];
    end
    if nargin < 8 || isempty(prefPitch)
        prefPitch = pi/2;
    end
    if nargin < 9 || isempty(family)
        family = 'up';
    end
    if nargin < 10 || isempty(fixedPitch)
        fixedPitch = false;
    end

    L = Links;
    D = L(6) + L(7);

    armJoints = zeros(1,6);
    isReachable = false;
    pitchUsed = NaN;

    % Plano horizontal
    r = hypot(x,y);
    if r <= L(5)
        return;
    end

    dFwd = sqrt(r^2-L(5)^2);
    hJ2 = z-L(1);

    % Restricoes
    ELBOW_MAX = deg2rad(140);
    PITCH_MAX = deg2rad(135);
    PITCH_MIN = deg2rad(30);
    STEP = deg2rad(5);

    if ~strcmp(family,'down')
        family = 'up';
    end

    if fixedPitch
        dList = 0;
    else
        dList = 0:STEP:(PITCH_MAX-PITCH_MIN);
    end

    % Procurar inclinacao valida
    for d = dList
        if d == 0
            cands = prefPitch;
        else
            cands = [prefPitch+d, prefPitch-d];
        end

        for psi = cands
            if ~fixedPitch && ...
                    (psi > PITCH_MAX || psi < PITCH_MIN)
                continue;
            end

            best = solve_pitch(psi,x,y,dFwd,hJ2, ...
                L,D,qMin,qMax,ELBOW_MAX,family);

            if isempty(best)
                continue;
            end

            % Refinamento por bissecao
            if d > 0
                good = psi;
                bad = psi-sign(psi-prefPitch)*STEP;

                for it = 1:15
                    mid = (good+bad)/2;
                    qm = solve_pitch(mid,x,y,dFwd,hJ2, ...
                        L,D,qMin,qMax,ELBOW_MAX,family);

                    if isempty(qm)
                        bad = mid;
                    else
                        good = mid;
                        best = qm;
                    end
                end
                psi = good;
            end

            % Orientacao da junta J6
            q6 = 0;
            if ~isempty(objYaw) && psi > deg2rad(45)
                q6 = mod(best(1)-objYaw+pi/4,pi/2)-pi/4;
            end

            armJoints = [best 0 q6];
            isReachable = true;
            pitchUsed = psi;
            return;
        end
    end
end
\end{lstlisting}

A função devolve os seis ângulos articulares, um indicador de existência de solução e a inclinação efetivamente utilizada. Quando não é encontrada uma configuração válida, o indicador \texttt{isReachable} permanece falso.

O argumento \texttt{fixedPitch} permite desativar a procura automática de inclinações alternativas, sendo utilizado, por exemplo, em movimentos verticais nos quais se pretende manter a mesma orientação da garra ao longo da trajetória.

As entradas e saídas da função encontram-se resumidas na Tabela~\ref{tab:inverse_io}.

\begin{table}[H]
    \centering
    \caption{Entradas e saídas da função \texttt{inverse\_kin}.}
    \label{tab:inverse_io}
    \small
    \begin{tabular}{lp{8.5cm}}
        \hline
        Argumento & Significado \\
        \hline
        \multicolumn{2}{l}{\textit{Entradas}} \\
        \texttt{x, y, z} & Posição pretendida no referencial da base [m] \\
        \texttt{Links} & Parâmetros geométricos do manipulador [m] \\
        \texttt{qMin, qMax} & Limites articulares [rad] \\
        \texttt{objYaw} & Orientação do objeto em torno de $z$, opcional [rad] \\
        \texttt{prefPitch} & Inclinação preferencial da garra, opcional [rad] \\
        \texttt{family} & Configuração \texttt{'up'} ou \texttt{'down'}, opcional \\
        \texttt{fixedPitch} & Utilização exclusiva da inclinação indicada, opcional \\
        \hline
        \multicolumn{2}{l}{\textit{Saídas}} \\
        \texttt{armJoints} & Vetor $[q_1\ q_2\ q_3\ q_4\ q_5\ q_6]$ [rad] \\
        \texttt{isReachable} & Indicador de existência de solução válida \\
        \texttt{pitchUsed} & Inclinação efetivamente utilizada [rad] \\
        \hline
    \end{tabular}
\end{table}

\subsection{Funções auxiliares da cinemática inversa}

A função \texttt{ik\_2dof} resolve a cinemática inversa do braço planar formado pelos elos $L_2$ e $L_3$, através da lei dos cossenos. A implementação permite selecionar a configuração com o cotovelo para cima ou para baixo, conforme descrito na Secção~\ref{sec:ik_lateral}.

Por sua vez, a função \texttt{solve\_pitch} recebe uma inclinação $\psi$ previamente definida, calcula a posição do pulso J4 e utiliza \texttt{ik\_2dof} para determinar os ângulos das juntas J2 e J3. Seguidamente, determina J1 e J4 e verifica se a configuração obtida respeita os limites articulares, a flexão máxima do cotovelo e as restrições de altura das juntas.

As implementações destas funções encontram-se nas Listagens~\ref{lst:ik2dof} e~\ref{lst:solvepitch}.

\begin{lstlisting}[
    language=Matlab,
    caption={Implementação da função ik\_2dof.},
    label={lst:ik2dof}
]
function [theta1, theta2, isReachable] = ...
    ik_2dof(xd, yd, LA, LB, solution_type)

    theta1 = NaN;
    theta2 = NaN;
    isReachable = false;

    r = sqrt(xd^2+yd^2);
    alpha1 = atan2(yd,xd);

    arg1 = (LA^2+r^2-LB^2)/(2*LA*r);
    arg2 = (LA^2+LB^2-r^2)/(2*LA*LB);

    if r == 0 || abs(arg1) > 1 || abs(arg2) > 1
        return;
    end

    beta1 = acos(arg1);
    beta2 = acos(arg2);

    if strcmp(solution_type,'down')
        theta1 = alpha1-beta1;
        theta2 = pi-beta2;
    else
        theta1 = alpha1+beta1;
        theta2 = -(pi-beta2);
    end

    isReachable = true;
end
\end{lstlisting}

\begin{lstlisting}[
    language=Matlab,
    caption={Implementação da função solve\_pitch.},
    label={lst:solvepitch}
]
function q = solve_pitch(psi,x,y,dFwd,hJ2, ...
    L,D,qMin,qMax,elbowMax,family)

    wrap = @(t) atan2(sin(t),cos(t));
    q = [];

    % Angulo da base
    q1 = atan2(y,x)-atan2(dFwd,L(5));

    % Posicao do pulso
    alpha3 = -psi;
    x3d = dFwd-D*cos(alpha3) ...
        -L(4)*cos(alpha3+pi/2);
    y3d = hJ2-D*sin(alpha3) ...
        -L(4)*sin(alpha3+pi/2);

    % Braco planar de dois elos
    [theta1,theta2,ok] = movement.ik_2dof( ...
        x3d,y3d,L(2),L(3),family);

    if ~ok
        return;
    end

    % Angulo da junta J4
    theta3 = (alpha3+pi/2)-theta1-theta2;

    % Angulos articulares
    q = wrap([q1,theta1-pi/2,theta2,theta3]);

    % Alturas das juntas
    zElbow = L(1)+L(2)*cos(q(2));
    zJ4 = zElbow+L(3)*cos(q(2)+q(3));
    zJ5 = zJ4+L(4)*cos(q(2)+q(3)+q(4));

    % Validacao da configuracao
    if any(q < qMin(1:4)) || ...
            any(q > qMax(1:4)) || ...
            abs(q(3)) > elbowMax || ...
            min([zElbow zJ4 zJ5]) <= 0
        q = [];
    end
end
\end{lstlisting}

As entradas e saídas das duas funções auxiliares encontram-se apresentadas nas Tabelas~\ref{tab:ik2dof_io} e~\ref{tab:solvepitch_io}.

\begin{table}[H]
    \centering
    \caption{Entradas e saídas da função \texttt{ik\_2dof}.}
    \label{tab:ik2dof_io}
    \small
    \begin{tabular}{lp{8cm}}
        \toprule
        Argumento & Significado \\
        \midrule
        \multicolumn{2}{l}{\textit{Entradas}} \\
        \texttt{xd, yd} & Coordenadas do pulso relativamente ao ombro [m] \\
        \texttt{LA, LB} & Comprimentos dos dois elos [m] \\
        \texttt{solution\_type} & Configuração \texttt{'up'} ou \texttt{'down'} \\
        \midrule
        \multicolumn{2}{l}{\textit{Saídas}} \\
        \texttt{theta1} & Orientação absoluta do primeiro elo [rad] \\
        \texttt{theta2} & Rotação do segundo elo relativamente ao primeiro [rad] \\
        \texttt{isReachable} & Indicador de solução geométrica válida \\
        \bottomrule
    \end{tabular}
\end{table}

\begin{table}[H]
    \centering
    \caption{Entradas e saídas da função \texttt{solve\_pitch}.}
    \label{tab:solvepitch_io}
    \small
    \begin{tabular}{lp{8cm}}
        \toprule
        Argumento & Significado \\
        \midrule
        \multicolumn{2}{l}{\textit{Entradas}} \\
        \texttt{psi} & Inclinação definida para a garra [rad] \\
        \texttt{x, y} & Coordenadas horizontais do alvo [m] \\
        \texttt{dFwd, hJ2} & Distância frontal e altura relativa ao ombro [m] \\
        \texttt{L, D} & Parâmetros geométricos do manipulador [m] \\
        \texttt{qMin, qMax} & Limites articulares [rad] \\
        \texttt{elbowMax} & Flexão máxima admitida para o cotovelo [rad] \\
        \texttt{family} & Configuração \texttt{'up'} ou \texttt{'down'} \\
        \midrule
        \multicolumn{2}{l}{\textit{Saída}} \\
        \texttt{q} & Vetor $[q_1\ q_2\ q_3\ q_4]$ [rad] ou vazio \\
        \bottomrule
    \end{tabular}
\end{table}

Deste modo, a implementação separa a procura de uma inclinação válida da resolução geométrica para uma inclinação específica, permitindo reutilizar os cálculos cinemáticos em diferentes procedimentos de movimento.

A função \texttt{inverse\_kin} é utilizada tanto no posicionamento do manipulador como no planeamento de trajetórias, sendo chamada repetidamente quando é necessário determinar configurações articulares para diferentes pontos do espaço de trabalho.
